\documentclass[10pt,a4paper]{amsart}
\usepackage[margin=19mm]{geometry}
\usepackage{amsmath,amssymb,mathtools}
\usepackage[scr=rsfs,cal=euler]{mathalfa}
\usepackage{xfrac}
\usepackage[unicode,hidelinks,bookmarks=false]{hyperref}
\newcommand{\ie}{i.e.\ }
\newcommand{\cf}{cf.\ }
\newcommand{\resp}{respectively\ }
\newcommand{\Subsection}{\subsection}
\providecommand{\nobreakdash}{\nobreak\hbox{-}}
\setlength{\emergencystretch}{2em}
\multlinegap0pt
\usepackage{bm}
\usepackage{bbm}
\usepackage{dsfont}
\usepackage{xspace}
\usepackage{cancel}
\usepackage{mathtools}
\usepackage{amsthm}
\makeatletter
\@ifundefined{lem}{\newtheorem{lem}{\bf Lemma}}{}
\makeatother
\newcommand{\mf}[1]{\mathbf{#1}}
\title[Asymptotic Expansion of the Gaussian Integral Operators on R]{Asymptotic Expansion of the Gaussian Integral Operators on Riemannian submanifolds of $\mathbb{R}^{n}$}
\author{Jia-Ming (Frank) Liou, Chi-Chien Lu}
\date{Source: arXiv:2506.13238v1 (CC BY 4.0)}
\begin{document}
\maketitle

\noindent\emph{Source and licence.} Asymptotic Expansion of the Gaussian Integral Operators on Riemannian submanifolds of $\mathbb{R}^{n}$. Version arXiv:2506.13238v1 (CC BY 4.0), distributed under the Creative Commons Attribution 4.0 International licence, as recorded in the arXiv licence metadata for this exact version and shown on the abstract page https://arxiv.org/abs/2506.13238v1. Full rights notice: licenses/arxiv-2506.13238-CC-BY-4.0.txt.\par\smallskip

\noindent\textit{Editorial scope.} Counted unit: the Lem whose source label is \texttt{tmetric} in the article (source0.tex lines 427--441, source label \texttt{tmetric}), delivered together with the article's own complete original proof of it (source0.tex lines 442--494, one \texttt{proof} environment of 53 lines). No mathematics is written, repaired or re-derived: statement and proof are the article's own text line for line. Required context retained uncounted: none. Every definition, display and label of those lines is retained unchanged. Omitted from this package and not redistributed: 1--426, 495--1354. Nothing inside a retained excerpt is omitted or altered: every retained body is delivered exactly as the article's lines are written, so no omission receipt of any kind accompanies this package. Citations: the retained excerpts carry no citation command at all, so no bibliography entry of the article is transcribed and no reference is redistributed. Reference adaptation: none was needed and none was made; every reference inside the delivered unit resolves inside it, so no unresolved reference or citation is produced. Printed slips kept literal: no printed word, symbol, hypothesis or number of the counted statement or of its proof was altered. Layout and class: the article's own class is \texttt{amsart} with packages including amsmath, amsfonts, graphics, epsfig, amssymb, amscd, xy, latexsym, graphicx, multirow, geometry, color; the wrapper is the lane's pinned \texttt{amsart} editorial wrapper at 10pt on A4, which restates the theorem-like environments the retained text needs and adds the article's own macro definitions that the retained text needs (\texttt{\textbackslash mf}), with extra wrapper packages (bm, bbm, dsfont, xspace, cancel, mathtools). The delivered numbering is the unit's own. Assets: no retained span contains a figure environment, a graphics-inclusion command, a PDF-page inclusion, a table or a TikZ picture; no image file is copied into this package, the package declares no asset dependency and no image file is redistributed with it. Licence: version arXiv:2506.13238v1 of this article is distributed under the Creative Commons Attribution 4.0 International licence, as recorded in the arXiv licence metadata for this exact version and shown on the abstract page https://arxiv.org/abs/2506.13238v1. A full rights notice is delivered at licenses/arxiv-2506.13238-CC-BY-4.0.txt. Characters: every retained excerpt is pure ASCII, so the delivered engine, XeLaTeX with its default Latin Modern text font, reports no missing character.
\begin{lem}\label{tmetric}
Let \( \mathbf{v} \in S^{d-1} \) be a unit vector. The function \( g_{\mathbf{v}} \) defined above admits a Taylor expansion around \( t = 0 \) of the form
\[
g_{\mathbf{v}}(t) = \sum_{j=0}^{\infty} \frac{g_j(\mathbf{v})}{j!} t^j, \quad \text{for } |t| < \delta.
\]
Each coefficient \( g_j(\mathbf{v}) \) is either zero or a homogeneous polynomial of degree \( j \) in \( \mathbf{v} = (v_1, \dots, v_d) \). The first few coefficients are given by
\[
g_0(\mathbf{v}) = 0, \quad 
g_1(\mathbf{v}) = 0, \quad 
g_2(\mathbf{v}) = 2, \quad 
g_3(\mathbf{v}) = 0, \quad 
g_4(\mathbf{v}) = -2 \| \mathbf{B}_x(\mathbf{v}_x, \mathbf{v}_x) \|^2,
\]
where \( \mathbf{B}_x(\mathbf{v}_x, \mathbf{v}_x) \) denotes the second fundamental form of the submanifold \( M \subset \mathbb{R}^n \) at the point \( x \), evaluated on the vector \( \mathbf{v}_x \in T_x M \).
\end{lem}

\begin{proof}
By definition, \( g_k(\mathbf{v}) = g_{\mathbf{v}}^{(k)}(0) \), so it suffices to compute the derivatives \( g_{\mathbf{v}}^{(k)}(0) \). Recall that
\[
g_{\mathbf{v}}(t) = \langle \gamma_{\mathbf{v}_{x}}(t) - x, \gamma_{\mathbf{v}_{x}}(t) - x \rangle_{\mathbb{R}^n}.
\]
Define \( \xi_{\mathbf{v}_{x}}(t) := \gamma_{\mathbf{v}_{x}}(t) - x \). Then \( \xi_{\mathbf{v}_{x}}(0) = 0 \), and \( \xi_{\mathbf{v}_{x}}^{(k)}(t) = \gamma_{\mathbf{v}_{x}}^{(k)}(t) \) for all \( k \geq 1 \). Since
\[
g_{\mathbf{v}_{x}}(t) = \langle \xi_{\mathbf{v}_{x}}(t), \xi_{\mathbf{v}_{x}}(t) \rangle,
\]
we apply the Leibniz rule to obtain
\[
g_{\mathbf{v}}^{(k)}(0) = \sum_{j=0}^{k} \binom{k}{j} \langle \xi_{\mathbf{v}_{x}}^{(j)}(0), \xi_{\mathbf{v}_{x}}^{(k-j)}(0) \rangle.
\]

The first few derivatives are:
\begin{align*}
g_{\mathbf{v}}'(0) &= 2 \langle \xi_{\mathbf{v}_{x}}(0), \xi_{\mathbf{v}_{x}}'(0) \rangle = 2 \langle 0, \mathbf{v}_{x} \rangle = 0, \\
g_{\mathbf{v}}''(0) &= 2 \| \gamma_{\mathbf{v}_{x}}'(0) \|^2 = 2 \| \mathbf{v}_{x} \|^2 = 2, \\
g_{\mathbf{v}}^{(3)}(0) &= 6 \langle \gamma_{\mathbf{v}_{x}}^{(2)}(0), \gamma_{\mathbf{v}_{x}}'(0) \rangle = 6 \langle \gamma_{\mathbf{v}_{x}}^{(2)}(0), \mathbf{v}_{x} \rangle.
\end{align*}

Since \( \gamma_{\mathbf{v}_{x}} \) is a geodesic in \( M \), it satisfies
\[
\langle \gamma_{\mathbf{v}_{x}}^{(2)}(t), \gamma_{\mathbf{v}_{x}}'(t) \rangle = 0 \quad \text{for all } t,
\]
so in particular,
\[
g_{\mathbf{v}}^{(3)}(0) = 6 \langle \gamma_{\mathbf{v}_{x}}^{(2)}(0), \mathbf{v}_{x} \rangle = 0.
\]

Differentiating the identity above with respect to \( t \), we obtain
\[
\langle \gamma_{\mathbf{v}_{x}}^{(3)}(t), \gamma_{\mathbf{v}_{x}}'(t) \rangle + \| \gamma_{\mathbf{v}_{x}}^{(2)}(t) \|^2 = 0.
\]
Evaluating at \( t = 0 \), we find
\[
\langle \gamma_{\mathbf{v}_{x}}^{(3)}(0), \mathbf{v}_{x} \rangle = - \| \gamma_{\mathbf{v}_{x}}^{(2)}(0) \|^2.
\]

Now we compute the fourth derivative:
\begin{align*}
g_{\mathbf{v}}^{(4)}(0)
&= 8 \langle \gamma_{\mathbf{v}_{x}}^{(3)}(0), \gamma_{\mathbf{v}_{x}}'(0) \rangle + 6 \| \gamma_{\mathbf{v}_{x}}^{(2)}(0) \|^2 \\
&= 8 \cdot \left( - \| \gamma_{\mathbf{v}_{x}}^{(2)}(0) \|^2 \right) + 6 \| \gamma_{\mathbf{v}_{x}}^{(2)}(0) \|^2 \\
&= -2 \| \gamma_{\mathbf{v}_{x}}^{(2)}(0) \|^2.
\end{align*}

Finally, using the fact that \( \gamma_{\mathbf{v}_{x}}^{(2)}(0) =\mf B_x(\mathbf{v}_{x}, \mathbf{v}_{x}) \), we obtain
\[
g_{\mathbf{v}}^{(4)}(0) = -2 \|\mf B_x(\mathbf{v}_{x}, \mathbf{v}_{x}) \|^2,
\]
as claimed.
\end{proof}

\end{document}
